\documentclass[10pt]{article}
\usepackage[margin=0.72in]{geometry}
\usepackage[T1]{fontenc}
\usepackage{lmodern}
\usepackage{amsmath,amssymb,amsthm}
\usepackage{booktabs}
\usepackage{microtype}
\usepackage[hidelinks]{hyperref}
\newtheorem{theorem}{Theorem}
\newcommand{\Q}{Q}
\newcommand{\Phiop}{\Phi}
\title{An All-Subgroup Containment-Threshold Atlas for a Frozen H\'enon Core}
\author{Anonymous}
\date{}

\begin{document}
\maketitle

\begin{abstract}
We replace exact full-core repair by a target-containment observable on the
twenty actual subgroup rows of
$\Q=\mathbb Z/9\oplus\mathbb Z/3\oplus\mathbb Z/2$.  For a deletion set
$D$ and retained labels $A=L\setminus D$, define
$\tau_H(D)=\min\{|R|:R\subseteq D,,H\subseteq\Phiop(A\cup R)\}$.
An exhaustive receipt stores all $20\times65536$ integer values, their
deleted-cardinality tables, and target distributions.  The full target row
is exactly the C78 repair distance, with distribution
$30400,32704,2368,64$.  An independent target-antichain reconstruction,
symbolic marginal checks, clean replay, and thirteen hostile mutations pass.
The result is a finite named-support atlas and makes no arithmetic, local,
Burnside-ring, table-of-marks, or Hilbert--Polya claim.
\end{abstract}

\section{Targeted repair}
Let $L=\{S_1,\ldots,S_{16}\}$ be the frozen named presentation.  For every
subgroup row $H\leq\Q$ and deletion set $D\subseteq L$, put
\[
 \tau_H(D)=\min\{|R|:R\subseteq D,\ H\subseteq
 \Phiop((L\setminus D)\cup R)\}.
\]
This is deliberately a containment threshold.  It does not require the
generated subgroup to equal $H$; in particular, different rows of the same
order remain separate named targets.

\section{Finite antichain identity}
\begin{theorem}
For a fixed target $H$, let $\mathcal M_H$ be the inclusion-minimal named
supports $M$ satisfying $H\subseteq\Phiop(M)$.  Then
\[
  \tau_H(D)=\min_{M\in\mathcal M_H}|M\setminus(L\setminus D)|.
\]
\end{theorem}
\begin{proof}[Proof sketch]
Any feasible restoration contains a target-containing support.  Removing
labels one at a time while containment persists reaches a member of
$\mathcal M_H$, and removing labels from the retained part is impossible by
the definition of the set difference.  Conversely, restoring $M\setminus A$
is feasible for every $M\in\mathcal M_H$.
\end{proof}

The C80 checker reconstructs the point-set addition table from C75, derives
each $\mathcal M_H$, and evaluates this formula for every one of the 65536
retained masks.  The twenty target orders in row order are
\[
1,2,3,3,3,3,6,6,6,6,9,9,9,9,18,18,18,18,27,54.
\]

\section{The exact-core boundary}
The last row is $H=\Q$.  Its threshold values are
\[
\#\{D:\tau_\Q(D)=0,1,2,3\}
  =(30400,32704,2368,64).
\]
The source-bound pivot/block decomposition gives, independently,
\[
 \tau_\Q(D)=\mathbf1_{S_9\in D}+
 \max\bigl(0,t(D)-2\bigr),
\]
where $t(D)$ counts fully deleted blocks of sizes $1,1,2,5$.  Thus the C80
row is exactly C78's $\rho$ row, while the other nineteen rows describe
proper-subgroup containment and can have smaller thresholds.

Every deleted-cardinality table has row sum $\binom{16}{k}$, hence each
target's cardinality marginal is $(1+x)^{16}$.  If $H_1\subseteq H_2$, then
\[
 \tau_{H_1}(D)\leq\tau_{H_2}(D)\quad\text{for every }D,
\]
which the receipt checks on all profiles using the actual subgroup bitsets.

\section{Certification and scope}
The producer binds the C75, C76, and C78 evidence/manifests by raw-byte
SHA-256.  The canonical C80 evidence hash is
\texttt{8d27428b14dbd7354e9c8308ad76b1108e3f551702165833301509cd52de7df5}.
The independent checker uses target antichains rather than the producer's
descending recurrence; a SymPy script verifies all twenty cardinality
marginals and the exact C78 polynomial; clean replay preserves the digest;
and $13/13$ hostile mutations are rejected.

This receipt concerns only the named sixteen-label finite model.  It does not
claim a full Burnside ring or table of marks, arithmetic/local data, Euler
factors, root numbers, automorphy, or a Hilbert--Polya operator.

\end{document}
